\section{Many-electron continuum estimate}\label{sec:many-body}

The point nuclei have now been fixed. We must compare their electronic
energy with the finite Hamiltonian $H_d$ on \emph{all} continuum states.
First we eliminate the space in which some electron lies outside the
localized modes, using the fixed gap and the two replacement estimates.
Then we identify the compressed Hamiltonian, keeping direct Coulomb
repulsion exactly at its leading order.

In the dilated coordinates the full operator after dividing by \(\lambda^2\) and shifting by \(-m(C_H+e_*)\) is
\[
             \mathcal H=\sum_{\ell=1}^m (L_p)_\ell+\lambda^{-1}\mathcal R,
           \qquad \mathcal R=\sum_{\ell<s}|y_\ell-y_s|^{-1}.
\]
We work on the antisymmetric space with spins and its $H^1$ form domain.
The Coulomb forms used here have the variational meaning specified in
Section~\ref{sec:introduction}, whether or not their energy infima are
attained. Indeed, fiberwise Hardy and Cauchy--Schwarz give, for every
$\eta>0$ and every fixed center $b$,
\[
 \int\frac{|u|^2}{|y_\ell-b|}
 \le 2\|u\|_2\|\nabla_\ell u\|_2
 \le \eta\|\nabla_\ell u\|_2^2+\eta^{-1}\|u\|_2^2.
\]
The same estimate applies to a pair singularity by fixing the other
coordinates. For each finite nuclear configuration, summing with
sufficiently small $\eta$ gives a closed lower-bounded form whose shifted
form norm is equivalent to the $H^1$ norm. Smooth compactly supported
antisymmetric functions form a core. The spectral infimum is therefore
the infimum of the form over normalized $H^1$ states, and the form
estimates below extend from this core by continuity.

Let $\Pi$ project onto the space in which all $m$ particles lie in $P$,
and write
\[
 \|u\|_1^2=\|u\|^2+\sum_\ell\|\nabla_\ell u\|^2.
\]
The smooth decaying modes lie in the form domain, so their compression is
well defined. The one-particle kinetic bound gives
$\|w\|_1\le p(n,D)\|w\|$ for $w\in\operatorname{Ran}\Pi$.

Let
\[
 E_\Pi=\min\operatorname{spec}
       \bigl((\Pi\mathcal H\Pi)|_{\operatorname{Ran}\Pi}\bigr),
\]
where the compression is defined by the form. This is a finite-dimensional
minimum because $\Pi$ projects onto the $m$-electron space built from
$2m$ one-particle spin modes.

\subsection{Control outside the orbital space}\label{subsec:complement-control}

\begin{proposition}[Eliminating the continuum complement]\label{prop:complement}
For the unit-nucleus construction above with $N=m$, and sufficiently
large polynomial scale $\lambda$, the full antisymmetric $m$-electron
operator in dilated and shifted units satisfies
\[
                 E_\Pi-p(n,D)\lambda^{-4/3}
                   \ \le\ \inf\operatorname{Spec}\mathcal H\ \le\ E_\Pi.       \tag{8}\label{eq:8}
\]
The lower bound holds on every normalized antisymmetric $H^1$ state,
including every spin sector.
\end{proposition}

\begin{proof}
The proof needs a lower bound on the complementary space, a bound on
its coupling to $\operatorname{Ran}\Pi$, and then a square completion.
Write $\mathcal H_c^0=\sum_\ell (L_c)_\ell$. By~\eqref{eq:6}, as
one-particle forms,
\[
                L_c\ge c'(I-P)-p(n,D)\lambda^{-0.9} I
\]
for some $c'>0$: the diagonal and cross terms involving $P$ are
controlled by $\|L_cP\|$. Since
$\sum_\ell(I-P)_\ell\ge I-\Pi$, summing gives a fixed positive lower
bound for $\mathcal H_c^0$ on $\operatorname{Ran}(I-\Pi)$ once
$m p(n,D)\lambda^{-0.9}$ is small. We also have the kinetic lower bound
$\mathcal H_c^0\ge-\sum_\ell\Delta_\ell/2-Cm$. A convex combination,
with a sufficiently small weight of order $1/m$ on the kinetic bound,
retains a positive $L^2$ term and gives
\[
                 \langle q,\mathcal H_c^0 q\rangle
                    \ge (c''/m)\|q\|_1^2,\qquad c''>0,
\]
for $q\in\operatorname{Ran}(I-\Pi)$. By~\eqref{eq:7} on each
coordinate and positivity of $\mathcal R$, this estimate persists with
half that coefficient for $\mathcal H$.

For $w\in\operatorname{Ran}\Pi$, equation~\eqref{eq:6} gives
$\|\mathcal H_c^0w\|\le p(n,D)\lambda^{-0.9}\|w\|$.
Hardy applied with one particle coordinate at a time also gives
$\|\mathcal Rw\|\le p(n,D)\|w\|$. Finally, polarizing the form
bound in~\eqref{eq:7} and using the $H^1$ bound on $w$ controls the
point-nucleus replacement in the mixed block. Together these estimates
give
\[
               |\langle q,\mathcal H w\rangle|
                  \le p(n,D)\lambda^{-2/3}\|w\|\|q\|_1.
\]
The same estimates show that the compressed form has norm less than
$c''/(4m)$ at sufficiently large scale. Set
$\alpha=c''/(4m)$ and let $\beta=p(n,D)\lambda^{-2/3}$ bound the
mixed block. Decompose any normalized form-domain vector as $u=w+q$,
where $w=\Pi u$ and $q=(I-\Pi)u$. Subtracting $E_\Pi$ makes the
$w$-diagonal form nonnegative and leaves at least
$\alpha\|q\|_1^2$ on the complement. Hence, with all pairings
interpreted as forms,
\[
 \begin{split}
 \langle u,\mathcal H u\rangle-E_\Pi\|u\|^2
 &\ge \alpha\|q\|_1^2-2\beta\|w\|\|q\|_1\\
 &\ge -\frac{\beta^2}{\alpha}\|w\|^2
 \ge -p(n,D)\lambda^{-4/3}.
 \end{split}
\]
This proves the lower bound in~\eqref{eq:8}; the upper bound follows by
testing a minimizing vector in $\operatorname{Ran}\Pi$. The lower-bound
argument applies to every normalized antisymmetric state, in any spin
sector, without any binding or ground-eigenvector assumption.
\end{proof}

\subsection{Compressed energy}\label{subsec:compressed-energy}

It remains to evaluate $\lambda E_\Pi$. Section~\ref{sec:density}
computed the one-particle matrix $\mathsf M$ in the orthonormal modes
$\widehat\psi_i$, to accuracy $p(n,D)\lambda^{-0.1}$. Multiplying the
compressed replacement error in~\eqref{eq:7} by $\lambda$ costs only
$p(n,D)\lambda^{-1/3}$, so the same matrix accuracy holds for the
point nuclei.
For the repulsion $\mathcal R$, the pair matrix on the original modes
before Gram correction has spatial entries from products of two
transition densities. Keeping only $\psi_i^2,\psi_j^2$ as the
respective densities gives the direct coefficients $v_{ij}$.
Any unequal indices within a density give an exponentially small error:
$\|\psi_i\psi_k\|_{L^1}=O(\lambda^{-0.8})$ for $i\ne k$, and the
Coulomb field of the other density product in absolute value is
uniformly bounded by orbital boundedness and decay. The matrix of
direct interactions in the product basis is bounded, and Gram
correction changes its entries negligibly.

These entrywise estimates give a many-electron operator-norm error of
at most $p(n,D)\lambda^{-0.1}$. To see the dimension dependence,
write the one- and two-body matrices in the $2m$ orthonormal spin modes.
Each coefficient multiplies a monomial $c_a^\dagger c_b$ or
$c_a^\dagger c_b^\dagger c_d c_c$, of norm at most one. There are
only $O(m^2)$ and $O(m^4)$ such monomials, respectively, so summing their
coefficient errors costs a polynomial factor on the entire fermion
space. In this representation, the retained two-body term is
\[
       \frac{U}{2}\sum_i n_i(n_i-1)+\sum_{i<j}v_{ij} n_i n_j .
\]
The retained one-body diagonal is
$-\sum_i(\sum_{j\ne i}v_{ij})n_i$. At the filling
$\sum_i n_i=m$, the identity~\eqref{eq:occupancy-completion} combines
these terms into the first line of~\eqref{eq:2}, including its scalar.
The hopping matrix gives the remaining line of the calibrated $H_d$.
We have thus identified the compressed Hamiltonian, and the comparison
with the full continuum follows.

\begin{proposition}[Continuum-to-fermion energy comparison]\label{prop:energy-comparison}
For the constructed unit nuclei with $N=m$, let $\mathcal H$ be the
dilated and shifted continuum operator above and let $H_d$ be the
finite fermion Hamiltonian~\eqref{eq:2}. At sufficiently large polynomial
scale,
\[
          |\lambda\inf\operatorname{Spec}\mathcal H-E(H_d)|
                              \le p(n,D)\lambda^{-0.1}.             \tag{9}\label{eq:9}
\]
\end{proposition}

\begin{proof}
The one- and two-body matrix estimates just obtained compare
$\lambda E_\Pi$ with $E(H_d)$ within $p(n,D)\lambda^{-0.1}$.
Multiplying~\eqref{eq:8} by $\lambda$ contributes only
$p(n,D)\lambda^{-1/3}$, which fits within that bound.
\end{proof}
This analysis treats long-range diagonal effects and possible charge transfer at the virtual-excitation scale in \eqref{eq:2}-\eqref{eq:3}, rather than requiring all intersite Coulomb matrix elements to be negligible compared to the input promise gap.

Table~\ref{tab:error-scales} records the different roles of the errors.
In particular the arbitrary-state form error cannot be used as a direct
energy error after multiplication by $\lambda$.

\begin{table}[tbp]
\centering
\small
\begin{tabular}{@{}p{.43\linewidth}p{.18\linewidth}p{.30\linewidth}@{}}
\toprule
Estimate & Error scale & Role in the comparison \\
\midrule
General form replacement, \eqref{eq:7}
 & $\lambda^{-2/3}$ & Controls mixed blocks in the $H^1$ norm \\
Compressed replacement, \eqref{eq:7}
 & $\lambda^{-4/3}$ & Direct error on the mode space \\
Complement elimination, \eqref{eq:8}
 & $\lambda^{-4/3}$ & Square of the mixed-block scale \\
Last two errors after multiplication by $\lambda$
 & $\lambda^{-1/3}$ & Error in the units of $H_d$ \\
Total comparison, \eqref{eq:9}
 & $\lambda^{-0.1}$ & Includes the other mode-matrix errors \\
\bottomrule
\end{tabular}
\caption{Error scales up to polynomial factors $p(n,D)$. The first three
rows use the units of $\mathcal H$; the last two use the units of $H_d$.
The $O(D^{-2})$ and $O(S^{-2})$ localization losses establish a fixed
complementary gap and do not enter the energy-precision budget.}
\label{tab:error-scales}
\end{table}

\subsection{Physical energies and thresholds}\label{subsec:physical-thresholds}

Choose the polynomial scale in $\lambda$ large enough that~\eqref{eq:9}
costs at most $\tau^2\gamma/100$ and that $\lambda\tau^2\gamma\ge8$.
All earlier smallness, separation, and calibration estimates hold
simultaneously by the stated growth choices. Undoing the dilation and
scalar shift gives the physical energy exactly as
\[
 E_0=\lambda^2\left[m(C_H+e_*)+
                         \inf\operatorname{Spec}\mathcal H\right].
\]
Combining~\eqref{eq:1}, \eqref{eq:3}, and~\eqref{eq:9} now gives a
total error at most $7\lambda\tau^2\gamma/100$, and hence at most
$\lambda\tau^2\gamma/8$, from
\[
       \mathcal F(E(H_{\rm in})),\qquad
       \mathcal F(\xi)=
           \lambda^2 m(C_H+e_*)+
           \lambda\left[-\sum_{i<j}v_{ij}
                    +\tau^2\left(\xi-C_s-\sum_e K_e\right)\right].            \tag{10}\label{eq:10}
\]
Take output thresholds rationally approximating \(\mathcal F(a_s+\gamma/4)\) and \(\mathcal F(b_s-\gamma/4)\) each within \(\lambda\tau^2\gamma/32\). For a YES source instance,
\[
 E_0\le\mathcal F(a_s)+\lambda\tau^2\gamma/8
       <\mathcal F(a_s+\gamma/4)-\lambda\tau^2\gamma/32\le a.
\]
For a NO source instance, the reversed estimate gives
\[
 E_0\ge\mathcal F(b_s)-\lambda\tau^2\gamma/8
       >\mathcal F(b_s-\gamma/4)+\lambda\tau^2\gamma/32\ge b.
\]
Since $b_s-a_s\ge\gamma$, the threshold separation satisfies
\[
 b-a\ge\lambda\tau^2\bigl(b_s-a_s-\gamma/2\bigr)
                   -\lambda\tau^2\gamma/16
      \ge\frac7{16}\lambda\tau^2\gamma\ge\frac72>1.
\]
Thus both promises survive rational rounding, and the reduction meets
the required output separation \(b-a\ge1\).
